\section{Representation size on stars}\label{sec:representation}
This appendix proves the lower bound on stars that
Section~\ref{sec:related} describes for the inverse-principal polytope
of Wei et al.\ \cite{WeiHull}. Remark~\ref{rem:dd-star} checks the
corresponding observation of Choi et al.\
\cite[Section~7.2]{ChoiDiagrams} for their decision diagrams.

Formulation size is measured with extended formulations, because
auxiliary variables can make a formulation much smaller. An
\emph{extended formulation}, or \emph{lift}, describes a convex set as
the projection of a higher-dimensional convex set that is easier to
write down. For example, the unit ball of the $\ell_1$ norm in $\R^n$
has $2^n$ facets, but it is the projection onto $x$ of
$\{(x,y):-y\le x\le y,\ \mathbf1^Ty\le1\}$, which needs only $2n+1$
inequalities. The size of the smallest lift, called the extension
complexity, therefore measures how hard a polytope is to formulate.

The basic example of a polytope without small lifts is the
\emph{correlation polytope}
$\operatorname{COR}(m)=\operatorname{conv}\{aa^T:a\in\{0,1\}^m\}$, also
known as the Boolean quadric polytope \cite[Section~3.2]{Fiorini}. Every
linear lift of $\operatorname{COR}(m)$ has $2^{\Omega(m)}$ inequalities
\cite[Theorem~7]{Fiorini}; the proof of Corollary~\ref{cor:conic}
recalls the corresponding bounds for second-order-cone and semidefinite
lifts. These bounds pass to every polytope with a face affinely
isomorphic to $\operatorname{COR}(m)$, because restricting a lift of the
polytope to the hyperplane that cuts out the face gives a lift of the
face that is no larger. Fiorini et al.\ derive their bounds for
stable-set and travelling-salesman polytopes in the same way, from faces
that project onto $\operatorname{COR}(m)$
\cite[Lemmas~8, 9, and~11]{Fiorini}.

Formally, a lift, or conic extended formulation, of a polytope $P$ over
a closed convex cone $K$ writes $P=\pi(K\cap\mathcal A)$ for an affine
subspace $\mathcal A$ and an affine map $\pi$; its size is measured by
the cone $K$. For a linear lift, $K$ is a nonnegative orthant $\R^r_+$,
and $r$ counts the inequalities. If an affine function $g$ is
nonnegative on $P$, then $\mathcal A'=\mathcal A\cap\{y:g(\pi(y))=0\}$
satisfies $\pi(K\cap\mathcal A')=P\cap\{g=0\}$, which gives a lift of
this face over the same cone. Composing with an affine isomorphism onto
$\operatorname{COR}(m)$ gives a lift of $\operatorname{COR}(m)$. Every
lower bound on lifts of the correlation polytope therefore applies to
any polytope with such a face.

For a positive-definite matrix $Q$ with index set $V$, the
\emph{inverse-principal polytope} is
\[
 P_Q=\operatorname{conv}\{(\mathbf1_S,W^S):S\subseteq V\},
 \qquad
 W^S_{SS}=Q_{SS}^{-1},\quad W^S_{ij}=0\text{ otherwise},
\]
with $W^\varnothing=0$. Each point $(\mathbf1_S,W^S)$, which we call an
atom, pairs the indicator vector of $S$ with the inverse of the
principal submatrix $Q_{SS}$, padded with zeros. This polytope yields a
convex-hull description of the epigraph of $x^TQx$ with indicators:
Wei et al.\ \cite[Section~3.1 and Theorem~1]{WeiHull} show that the
closed convex hull of
$\{(x,z,t):t\ge x^TQx,\ x_i(1-z_i)=0,\ z\in\{0,1\}^V\}$ is the set of
points for which some $W$ satisfies $(z,W)\in P_Q$ and
$\left(\begin{smallmatrix}W&x\\x^T&t\end{smallmatrix}\right)\succeq0$.
A small extended formulation of $P_Q$ therefore gives a small extended
formulation of this hull. Liu et al.\ \cite{LiuStieltjes} study
Stieltjes matrices, the positive-definite matrices with nonpositive
off-diagonal entries. They prove hardness of optimizing over a related
inverse-graph closure for a diagonal-minus-rank-one matrix
\cite[Section~4.3, Proposition~6]{LiuStieltjes}. We show that $P_Q$
itself has no small exact extended formulation, even for a sparse,
uniformly conditioned star.

\begin{theorem}\label{thm:star-face}
For $m\ge1$, let $b=1/(2m)$ and
\[
 Q=\begin{pmatrix}1&b\mathbf1^T\\b\mathbf1&I_{2m}\end{pmatrix},
\]
with rows and columns indexed by the root $0$ and the leaves
$1,\ldots,2m$. Then $Q$ has a star support graph, rational entries of
$O(\log(m+1))$ bits, and condition number below six. The polytope $P_Q$
has an exposed face affinely isomorphic to the correlation polytope
$\operatorname{COR}(m)=\operatorname{conv}\{aa^T:a\in\{0,1\}^m\}$.
The same holds for the Stieltjes matrix obtained by changing the sign of
the root coordinate.
\end{theorem}
The construction works because, when the root is selected, inverting
$Q_{SS}$ couples any two selected leaves through the root: the inverse
entry between distinct leaves $i$ and $j$ is then a multiple of
$z_iz_j$, like the entry $a_ia_j$ of $aa^T$. The multiple depends on how
many leaves are selected. To keep it fixed, the proof pairs leaf $i$
with leaf $m+i$ and uses an affine function that vanishes exactly at the
atoms selecting the root and one leaf of each pair. These atoms select
$m$ leaves, so on them the inverse entry between distinct leaves
$i,j\le m$ is a fixed multiple of $z_iz_j$, as in
$\operatorname{COR}(m)$.
\begin{proof}
The eigenvalues of $Q$ are one and $1\pm1/\sqrt{2m}$, so its condition
number is at most $(1+1/\sqrt2)/(1-1/\sqrt2)=3+2\sqrt2<6$.

We first compute the atoms. Write $L$ for the leaves, $z\in\{0,1\}^{2m}$
for the leaf indicators of an atom, and $k=\mathbf1^Tz$. If the root is
not selected, the leaf block of $W$ is $\operatorname{Diag}(z)$. If the
root is selected, put $\omega=(1-b^2k)^{-1}$, which is finite because
$b^2k\le1/(2m)<1$. Block inversion gives
\begin{equation}\label{eq:star-inverse}
 W_{00}=\omega,\qquad W_{0L}=-b\omega z^T,\qquad
 W_{LL}=\operatorname{Diag}(z)+b^2\omega zz^T.
\end{equation}

Next we expose the face. Consider the affine function
\[
 g(z_0,z,W)=(m+1)(1-z_0)+m-\mathbf1^Tz
                         +b^{-2}\sum_{i=1}^mW_{i,m+i}.
\]
At an atom without the root, $g=2m+1-k\ge1$. At an atom with the root,
let $p$ be the number of pairs $\{i,m+i\}$ with both leaves selected.
Then $b^{-2}\sum_iW_{i,m+i}=\omega p$ and $k\le m+p$, so
$g=m-k+\omega p\ge(\omega-1)p\ge0$. If $p>0$, then $k\ge2$ and
$\omega>1$, so $g>0$. Hence $g\ge0$ on $P_Q$, and an atom has $g=0$
exactly when the root and
exactly one leaf of each pair are selected, that is, when
$z=(a,\mathbf1-a)$ for some $a\in\{0,1\}^m$. The hyperplane $g=0$
therefore exposes the face $\mathcal F$ of $P_Q$ spanned by these atoms.

Finally we identify $\mathcal F$ with $\operatorname{COR}(m)$. On
$\mathcal F$, $z_0=1$ and $k=m$, so $\omega=(1-mb^2)^{-1}$ is
constant. Define the affine map $\phi$ by $X_{ii}=z_i$ and
$X_{ij}=W_{ij}/(b^2\omega)$ for distinct $i,j\in\{1,\ldots,m\}$. By
\eqref{eq:star-inverse}, $\phi$ sends the atom with $z=(a,\mathbf1-a)$
to $aa^T$, so $\phi(\mathcal F)=\operatorname{COR}(m)$. For the inverse,
let $a=\operatorname{diag}(X)$, $z=(a,\mathbf1-a)$, and
\[
 R(X)=\begin{pmatrix}
 X&a\mathbf1^T-X\\
 \mathbf1a^T-X&\mathbf1\mathbf1^T-a\mathbf1^T-\mathbf1a^T+X
 \end{pmatrix}.
\]
Let $\psi(X)$ be the point with $z_0=1$, leaf indicators $z$,
$W_{00}=\omega$, $W_{0L}=-b\omega z^T$, and
$W_{LL}=\operatorname{Diag}(z)+b^2\omega R(X)$. The map $\psi$ is affine,
and $R(aa^T)=zz^T$, so $\psi(aa^T)$ is the atom \eqref{eq:star-inverse}.
Thus $\psi\circ\phi$ is the identity at the atoms of $\mathcal F$, and
$\phi\circ\psi$ at the points $aa^T$. Both maps are affine, so they are
the identity on $\mathcal F$ and on $\operatorname{COR}(m)$.

For the sign change, let $D=\operatorname{Diag}(-1,1,\ldots,1)$. Every
padded principal inverse of $DQD$ equals $DW^SD$, so the invertible
linear map $(z,W)\mapsto(z,DWD)$ maps $P_Q$ onto $P_{DQD}$.
\end{proof}

\begin{corollary}\label{cor:conic}
Let $Q$ be either matrix in Theorem~\ref{thm:star-face}. An exact lift of
$P_Q$ needs exponentially many nonnegative scalar factors if it is a
linear program, exponential total cone dimension if it uses second-order
cones, and exponentially many positive-semidefinite blocks if every block
has order at most a fixed $d$. Scalar inequalities count as blocks of
order one. A semidefinite lift with blocks of arbitrary order needs total
block order $2^{\Omega(m^{2/13})}$.
\end{corollary}
\begin{proof}
By the restriction argument before Theorem~\ref{thm:star-face}, it
suffices to bound lifts of $\operatorname{COR}(m)$ over the same cones.
For fixed $d$, Fawzi and Parrilo \cite[Theorem~1]{FawziParrilo} show
that the number of blocks of order $d$ grows exponentially in $m$;
blocks of smaller order embed in blocks of order $d$. A linear program
is the case $d=1$, where the bound is due to Fiorini et al.\
\cite[Theorem~7]{Fiorini}. For $d=2$, at least $7^{-1/2}(9/7)^{m/2}$
blocks are needed. A Lorentz cone of dimension $k\ge3$ has a lift with
$k-2$ three-dimensional Lorentz cones (a standard construction, recalled
in \cite[p.~3]{FawziParrilo}), each linearly isomorphic to the
positive-semidefinite cone of order two. Lorentz cones of dimensions one
and two are linearly isomorphic to products of one or two scalar cones.
A second-order-cone lift of total dimension $D$ thus gives a lift with at
most $D$ blocks of order at most two. For blocks of arbitrary order,
place them in one block-diagonal matrix whose order is their total
order, and apply Lee, Raghavendra, and Steurer
\cite[Theorems~1.1 and 5.4]{LeeSDP}.
\end{proof}

The lower bound concerns $P_Q$ together with all its inverse entries.
The obstruction lies in the leaf--leaf entries of $W$, which belong to
pairs of vertices that are not adjacent in the star, and it disappears
when they are projected out. On $\mathcal F$, the indicators, the
diagonal of $W$, and the entries $W_{0i}$ on star edges are affine
functions of $a=(z_1,\ldots,z_m)$, so their projection is affinely
isomorphic to the cube $[0,1]^m$. The argument therefore gives no lower
bound for formulations of the indicator quadratic epigraph that
eliminate $W$.

\begin{remark}[Projected-row decision diagrams on stars]\label{rem:dd-star}
The decision diagrams of Choi et al.\ \cite{ChoiDiagrams} fix the
indicators one at a time in a chosen variable order. Each node of a
layer is a state that summarizes how the indicators fixed so far affect
the rest of the problem, and partial assignments with the same state
share a node; the size of a diagram is its number of nodes. Choi et al.\
\cite[Section~7.2 and Corollary~1]{ChoiDiagrams} bound the number of
nodes of their exact decision diagrams on a tree with $k$ leaves by a
polynomial of degree $k+1$ in the number of variables, for a variable
order in which every vertex of the tree precedes its parent. They note
that on stars the resulting representation becomes exponential in size,
although tree problems are solvable in polynomial time. Their states
store the future rows of a factor $F$ with $Q=FF^T$, projected
orthogonally away from the rows of the selected past variables
\cite[Definition~7 and Lemma~2]{ChoiDiagrams}; their Section~EC.2.2
also allows normalizing each stored row.

A direct check shows that these diagrams are exponential on stars for
every variable order. Let $Q$ have unit diagonal, root $0$, leaves
$1,\ldots,n$, entries $Q_{0i}=b_i\ne0$, and $\sum_ib_i^2<1$ (for
example, $b_i=1/(2n)$, so that $\|Q-I\|_2=1/(2\sqrt n)$). With
unconstrained supports, the layer just before the last leaf $j$ in the
order has at least $2^{n-1}$ distinct states, with or without row
normalization. Indeed, the rows $f_0,\ldots,f_n$ of $F$ are linearly
independent, the leaf rows are orthonormal, and
$\langle f_0,f_i\rangle=b_i$. Let $S$ be the set of selected leaves
other than $j$. If the root follows $j$, the state contains the root row
$r_0(S)=f_0-\sum_{i\in S}b_if_i$. If the root precedes $j$ and is
selected, the state is the row $r_j(S)=f_j-(b_j/d_S)\,r_0(S)$ of $j$,
where $d_S=\|r_0(S)\|^2=1-\sum_{i\in S}b_i^2>0$. In the basis
$f_0,\ldots,f_n$, each row has one coefficient equal to one, and its
coefficients on leaves $i\ne j$ are nonzero exactly for $i\in S$.
Different subsets thus give nonproportional rows, which row scaling
cannot merge. This confirms their observation and is not a contribution
of this paper.
\end{remark}
